\chapter{Espace vectoriel}
\section{Espace vectoriel sur un corps}


\lecture
\begin{definition}[Espace vectoriel] \label{def: Espace vectoriel}
Soit $\K$ un corps, un $\K$\textbf{-espace vectoriel} $V$ est groupe commutatif muni d'une loi de composition externe
	\[
	\begin{array}{cccc}
		\cdot * \cdot : & \K \times V & \to & V \\
						& (\lambda, v) & \mapsto & \lambda * v
	\end{array}
	\]
ayant les propriétés :
\begin{description}
	\item[Associativité :] $\forall \lambda_1, \lambda_2 \in \K$ et $v \in V$,
	\[
	(\lambda_1 \cdot \lambda_2) * v = \lambda_1 * (\lambda_2 * v)
	\]
	\item[Distributivité :] $\forall \lambda_1, \lambda_2 \in \K$ et $v_1, v_2 \in V$,
	\[
	(\lambda_1 + \lambda_2) * v_1 = \lambda_1 * v_1 + \lambda_2 * v_1 \et \lambda_1 * (v_1 + v_2) = \lambda_1 * v_1 + \lambda_1 * v_2
	\]
	\item[Neutralité de $1$ :] $\forall v \in V$,
	\[
	1 * v = v	
	\]
\end{description}
Les éléments $v$ de $V$ sont appelés des \textbf{vecteurs} tandis que les éléments $\lambda \in \K$ sont appelés des \textbf{scalaire}.
\end{definition}


\begin{definition}[$\K$-algèbre] \label{def: algèbre}
Soit $\K$ un corps, une $\K-$\textbf{algèbre} est un anneau $(B, +, \cdot)$ possédant une structure de $\K-$espace vectoriel qui vérifie la propriétés d'associativité suivant,
	\[
	\forall \lambda \in \K, ~\ v_1, v_2 \in B, \quad \lambda * (v_1 \cdot v_2) = (\lambda * v_1) \cdot v_2 =  v_1 \cdot (\lambda * v_2)
	\]
\end{definition}


\subsection{Sous-espace vectoriel}

\begin{definition}[Sous-espace vectoriel] \label{def : Sous-espace vectriel}
Soit $\K$ un corps et $V$ un $\K$-espace vectoriel. Un \text{sous-espace vectoriel} $W \subset V$ d'un $\K$-espace vectoriel est un sous-groupe de $(V, +)$ qui est stable par la multiplication par les scalaires,
	\[
	\forall \lambda \in \K, ~\ \forall w \in W, \quad \lambda * w \in W
	\]
\end{definition}

\begin{prop}[Critère de sous-espace vectoriel] \label{prop: Critère de sous-espace vectoriel}
Soit $W \subset V$ un sous-ensemble non-vide d'un $\K$-espace vectoriel, alors $W$ est une sous-espace vectoriel si et seulement si
	\begin{equation} \label{eq: eq1}
	\forall \lambda \in \K, ~\ \forall w_1, w_2 \in W, \quad \lambda w_1 + w_2 \in W
	\end{equation}
\end{prop}


\prf{Pour tout $w_1, w_2 \in W$, et en application la condition \eqref{eq: eq1} à $w_1, w_2$ et $\lambda = -1$ on a,
	\[
	w_1 - w_2 \in W
	\]
	donc $W$ vérifie le critère de sous-groupe et est donc un sous-groupe de $(V,+)$. Il contient en particulier $0_V$ et alors pour tout $\lambda \in \K$, on a
	\[
	\forall w \in W, \quad \lambda * w + 0_V \in W
	\]
}

\subsection{Sous-espace vectoriel engendré par un ensemble}
\begin{definition}[Sous-espace vectoriel engendré] \label{def : Sous-espace vectoriel engendré}
Soit $X \subset V$ un sous-ensemble d'un $\K$-espace vectoriel, le \textbf{module engendré} par $X$ est le petit petit sous-module de $V$ contenant $X$,
	\[
	\langle X \rangle_\K := \bigcap_{X \subset N \subset V} N
	\]
\end{definition}

\begin{prop}Soit $V$ un $\K$-espace vectoriel et $X \subset V$ un sous-ensemble de $V$ alors $\langle X \rangle_\K$ est soit le sous-espace vectoriel trivial $\{0_V\}$ si $X$ est vide, soit l'ensemble des combinaisons linéaires de $X$ à coefficients dans $\K$,
	\[
	\langle X \rangle_\K = \CL_\K(X) := \left\{ \sum_{i=1}^n \lambda_i * x_i \mid n \geqslant 1, ~ \lambda_i \in \K, ~ x_i \in X\right\}
	\]
\end{prop}

\prf{On suppose $X$ non-vide. Soit $N \supset X$ un sous-espace vectoriel contenant $X$ alors pour tout $\lambda_1, \ldots, \lambda_n \in \K$ et tout $x_1, \ldots x_n \in X$ on a
	\[
	\sum_{i=1}^n \lambda_i * x_i \in N
	\]
	par stabilité de $N$. Donc tout sous-espace vectoriel $N$ contenant $X$ contient $\CL_\K(X)$. 
	
	Il reste donc à montrer que $\CL_\K(X)$ est un sous-espace vectoriel, alors soit $v, w \in \CL_\K(X)$ deux combinaisons linéaire d'éléments de $X$, alors il est facile de voir que,
	\[
	\lambda * v + w \in \CL_\K(X)
	\]
	ainsi par le critère de sous-espace vectoriel $\CL_\K(X)$ \eqref{prop: Critère de sous-espace vectoriel} est bien un sous-espace vectoriel.
}


\begin{definition}[Famille génératrice] \label{def: famille génératrice}
	Si $\langle X \rangle_\K = V$, on dit que $X$ est une \textbf{famille génératrice} de $V$, de plus on dit que $V$ est de \textbf{type fini} si $X$ est fini.
\end{definition}

\subsection{Applications linéaires}
\begin{definition}[Application linéaire] \label{def: application linéaire}
	Soit $\K$ un corps et $V, W$ deux $\K$-espace vectoriel, une \textbf{application linéaire} entre $V$ et $W$ est un morphisme de groupe,
	\[
	\varphi : V \to W
	\]
	qui est compatible avec les lois de compositions externes
	\[
	\forall \lambda \in \K, ~\ \forall v \in V, \quad \varphi(\lambda * v) = \lambda * \varphi(v)
	\]
	On notera
	\[
	\Hom_\K(V,W) \et \End_\K(V)
	\]
	l'ensemble des applications linéaire $\varphi : V \to W$ et $\varphi : V \to V$ respectivement.
\end{definition}

\begin{prop}[Critère d'application linéaire] \label{prop: Critère d'application linéaire}
	Une application entre espaces vectoriels $\varphi$ est une application linéaire si et seulement si 
	\begin{equation}\label{eq: Critère d'application linéaire}
	\forall \lambda \in \K, ~\ v_1, v_2 \in V, \quad \varphi(\lambda * v_1 + v_2) = \lambda * \varphi (v_1) + \varphi(v_2)
	\end{equation}	
\end{prop}

\prf{On applique \eqref{eq: Critère d'application linéaire} avec $\lambda = 1$. On a donc
	\[
	\forall v_1, v_2 \in V, \quad \varphi(v_1 + v_2) = \varphi(v_1) + \varphi(v_2)
	\]
	donc $\varphi$ est un morphisme de groupe. On a donc $\phi(0_V) = 0_W$ et donc
	\[
	\forall \lambda \in \K, ~\forall v \in V, \quad \varphi(\lambda * v_1) = \varphi(\lambda * v_1 + 0_V) = \lambda * \varphi(v_1) + 0_W = \lambda + \varphi(v_1)
	\]
}

\begin{thm}[Espace des applications linéaires] \label{thm: Espace des applications linéaires}
L'ensemble des applications linéaires $\Hom_\K(V,W)$ possède un structure naturelle de $\K$-espace vectoriel.
\end{thm}

\prf{On sait qu'une combinaison de deux applications linéaires et linéaire, c'est-à-dire
	\[
	\forall \varphi, \phi \in \Hom_\K(V,W), ~\forall \lambda \in \K, \quad \lambda.\varphi + \phi : v \mapsto \lambda \varphi(v) + \phi(v)
	\]
	est $\K$-linéaire. Donc le théorème en découle facilement.
}

\lecture

\section{Famille génératrice, libre et base}
\begin{definition}[Famille génératrice] \label{def: famille génératrice espace vectoriel}
	Soit $V$ un $\K$-espace vectoriel. Un sous-ensemble $G \subset V$ est une \textbf{famille génératrice} si
	\[
	\Vect (G) = \langle G \rangle_\K = V
	\]
	Si $V$ admet une famille génératrice finie, on dit que $V$ est un $\K$-espace vectoriel fini.
\end{definition}

\begin{definition}[Dimension] \label{def: dimension d'un espace vectoriel}
	Soit $V$ un $\K$-espace vectoriel de type fini. Si $V$ est non-nul, sa dimension est le cardinal minimum d'une famille génératrice :
	\[
	\dim V := \min_{G \text{ génératrice}} |G|
	\]
	On dira alors que $V$ est de \textbf{dimension finie} au lieu de type fini.
\end{definition}
Supposons que $G = \{ e_1, \ldots, e_d\} \subset V$ soit une famille génératrice finie de $V$ de cardinal $d \geqslant \dim V$. Tout élément $v \in V$ peut donc se représenter sous la forme d'une combinaison linéaire des $e_i$,
\[
v = \sum_{i=1}^d v_i e_i, \quad v_i \in \K
\]
En d'autres termes, on dispose d'une application "combinaison linéaire" qui est surjective
\[
\begin{array}{cccc}
	\CL_G : & \K^d & \to & V \\
			& (v_1, \ldots, v_d) & \mapsto & \CL_G(v_1, \ldots, v_d) = v_1 e_1 + \ldots + v_d e_d
\end{array}
\]
\begin{lem}
	L'application $\CL_G$ est linéaire.
\end{lem}
\prf

